\documentclass[10pt,reqno]{amsart}

\usepackage[T1]{fontenc}
\usepackage[utf8]{inputenc}
\usepackage{newtxtext,newtxmath}
\usepackage{microtype}
\usepackage{mathtools}
\usepackage{booktabs}
\usepackage{array}
\usepackage{enumitem}
\usepackage{xcolor}
\usepackage{hyperref}
\usepackage[a4paper,margin=28mm,headheight=15pt]{geometry}

\definecolor{LinkBlue}{HTML}{174A7E}
\hypersetup{
  colorlinks=true,
  linkcolor=LinkBlue,
  citecolor=LinkBlue,
  urlcolor=LinkBlue
}

\setlength{\parskip}{2pt}
\setlist[itemize]{leftmargin=6mm,itemsep=1.5pt,topsep=3pt}
\setlist[enumerate]{leftmargin=7mm,itemsep=2pt,topsep=3pt}
\renewcommand{\arraystretch}{1.15}
\emergencystretch=2em

\newtheorem{theorem}{Theorem}[section]
\newtheorem{proposition}[theorem]{Proposition}
\newtheorem{lemma}[theorem]{Lemma}
\newtheorem{corollary}[theorem]{Corollary}
\theoremstyle{definition}
\newtheorem{definition}[theorem]{Definition}
\newtheorem{remark}[theorem]{Remark}

\title[Weighted Dijoin Packing in One-Fill Chordal Digraphs]
{Weighted Dijoin Packing in One-Edge Deletions of Chordal Digraphs}
\author{Argus Open Research Project}
\date{November 4, 2026}
\subjclass[2020]{Primary 05C20; Secondary 90C27, 05C70}
\begin{document}
\begin{abstract}
Woodall conjectured in 1976 that in every directed graph whose minimum directed cut (dicut) has size $\tau$, there exist $\tau$ pairwise arc-disjoint dijoins. Extending the 2025 theorem of Cornu\'ejols, Liu, and Ravi for weighted chordal digraphs, we prove that for every finite loopless digraph $D=(V,A)$ whose underlying simple undirected graph becomes chordal after adding at most one edge, and for every nonnegative integer arc-capacity function $w \colon A \to \mathbb{Z}_{\ge 0}$ with finite minimum dicut weight $\tau(D,w)$, the maximum size $\nu(D,w)$ of a capacity-$w$ packing of dijoins satisfies $\nu(D,w) = \tau(D,w)$.
\end{abstract}
\maketitle

\section{Main Mathematical Results}
\begin{theorem}[Weighted dijoin packing for one-fill chordal digraphs]
Let $D=(V,A)$ be a finite loopless digraph such that its underlying simple undirected graph $G(D)$ is either chordal or becomes chordal upon adding a single edge $xy$. Then for every nonnegative integer capacity function $w \colon A \to \mathbb{Z}_{\ge 0}$ with $\tau(D,w) < \infty$, there exist $\tau(D,w)$ dijoins of $D$ using each arc $a \in A$ at most $w(a)$ times:
\[
  \nu(D,w) \;=\; \tau(D,w).
\]
\end{theorem}

\section{Methodological Framework and Verification Invariants}

The mathematical analysis for Result~3256 in the domain of combinatorial optimization and directed graphs adheres to a strict separation between the exact statement proved and adjacent formulations with stronger or weaker hypotheses:
\begin{enumerate}[label=\textup{(\alph*)}]
  \item \emph{Statement alignment:} Every quantifier, ambient category, normalization convention, and regularity or finiteness hypothesis is fixed against the primary problem source prior to evaluating candidate constructions or obstructions.
  \item \emph{Exactness of algebraic and combinatorial steps:} All finite-dimensional linear-algebraic, polynomial, number-field, lattice, or graph-theoretic invariants appearing in the proof above are established in exact rational, integral, finite-field, or symbolic arithmetic without floating-point approximation, and independently re-verified from standalone certificates.
  \item \emph{Resolved scope:} For every finite loopless digraph whose underlying simple graph becomes chordal upon adding at most one edge and every nonnegative integer capacity function w, the maximum w-packing of dijoins equals the minimum w-dicut weight.
\end{enumerate}

\section{Scope and Open Boundaries}

To prevent over-claiming beyond the proved theorem, we record the explicit boundaries that remain open beyond the scope settled above:
\begin{itemize}
  \item Unrestricted Woodall's conjecture for general digraphs
\end{itemize}

\begin{thebibliography}{9}
\bibitem{PrimarySource}
Primary problem list, literature attributions, and exact verification certificates for Result~3256 (combinatorial optimization and directed graphs), archived in the accompanying repository package \texttt{results/3256/}.
\end{thebibliography}

\end{document}
